\section{Introduction}\label{sec:intro}
\paragraph*{Equational Theories of Regular Languages.}
One of the landmark achievements in the theory of automata is the axiomatisation of Kleene algebra. In his original paper on regular expressions~\cite[p.~32]{Kleene1956}, Kleene asked: ``Is there an effective procedure for deciding whether two regular expressions are equal, or equivalent?'' A positive answer to this question is now  known to be a simple consequence of decidable equivalence for finite automata, however the question was quickly picked up in the following meaning: is there an equational axiomatisation, in which every two equivalent regular expressions can be transformed into each other by using a sequence of valid equations such as  $(L +K)^* = L^* \cdot (K  \cdot L^*)^*$? 

In 1964, Redko proved that there cannot be any  axiomatisation that consists of a finite set of equations~\cite{Redko1964}. However, the problem continued to attract a lot of attention, including an entire book by Conway~\cite{Conway1971}. One line of work keeps the axiomatisation purely equational, at the cost of using infinitely many equations. Conway proposed such an infinite family of equational axioms, containing certain \emph{monoid identities}~\cite[pp.~116--119]{Conway1971}, and its completeness was later proved by Krob~\cite[Theorem 0.1]{Krob1991}. Around the same time, Bloom and \'Esik proved completeness for another infinite family of equations, in which the monoid identities are replaced by more general \emph{commutative identities}~\cite[Corollary 2.14]{Bloom_Ésik_1993}. Another line of work strives for axiomatisations that are finite, but at the cost of using certain rules instead of only equations. An important example is Salomaa's axiomatisation~\cite[Theorems 1 and 2]{Salomaa1966}, which uses certain induction rules. Another example -- which is the state of the art in the field -- is the axiomatisation by Kozen~\cite[Theorem 5.5]{Kozen1994}, which uses a finite set of equations and two implications: 
\begin{align*}
L \cdot K \subseteq K \quad \text{implies} \quad L^* \cdot K \subseteq K
\qquad \text{and} \qquad 
K \cdot L \subseteq K \quad \text{implies} \quad K \cdot L^* \subseteq K.
\end{align*}
Subsequent work gave complete axiomatisations of many related systems, including: Kleene algebra without the identity element~\cite[Theorems 10 and 11]{DoumanePous2018}, Kleene algebra with tests~\cite[Theorem 8]{KozenSmith1996}, a network version of Kleene algebra with tests~\cite[Section 4]{Anderson2014NetKAT}, a version of Kleene  
algebra for tree languages~\cite[Theorem 8]{Esik1998}, and a version of Kleene algebra for regular expressions quotiented by bisimulation~\cite[Theorem 5.1]{Grabmayer2022}. Other relevant work includes complete axiomatisations of: iteration theories~\cite[Chapter 14]{BloomEsik1993}, algebras that describe graphs of treewidth at most three~\cite[Theorem 5.5]{DoumaneHumeauPous2024}, and the propositional $\mu$-calculus~\cite[Theorem 37]{Walukiewicz2000}.

\paragraph*{Transducers.}
The work on Kleene algebra concerns regular languages, which can be seen as string-to-Boolean functions. The goal of this paper is to study the equational theories of transducers,  which define string-to-string functions. To the best of our knowledge, this is the first work of this kind.
The particular transducer model that is studied in this paper is the \emph{rational functions}, which we describe below, although in future work we intend to consider more powerful models, such as the regular functions~\cite[Part C]{Bojanczyk2026}, or even the polyregular functions~\cite{Bojanczyk2022}.

A good way to describe the \kl{rational functions}, which highlights the connections with the classical work on regular expressions, is to think of them as a special case of the more general \emph{rational relations}. Rational relations are defined to be binary relations between strings that can be defined by a regular expression in the product monoid $A^* \times B^*$, and rational functions are the \kl{rational relations} that are also functions. They are one of the central models in the transducer literature, with numerous equivalent characterisations, including:  compositions of sequential machines~\cite[Proposition 7.4]{Elgot1965}, Eilenberg  bimachines~\cite[Chapter XI]{Eilenberg1974}, relabelings definable in \mso~\cite[Section C.4.2]{Bojanczyk2026}, as well as origin transductions with monotone origins~\cite[Theorem 3]{Bojanczyk2014}.

A natural idea would be to find a complete axiomatisation of the rational relations. However, this is not possible, since such an axiomatisation (at least with a recursively enumerable set of axioms) would imply that equivalence for rational relations is recursively enumerable, which is known to be false (see~\cite[p.~412]{Griffiths1968} and~\cite[Theorem 7]{FischerRosenberg1968}). 
This problem does not occur for rational functions as their equivalence is decidable, which can be proved using a canonical form shown in~\cite[Theorem 2]{ReutenauerSchutzenberger1991}.\footnote{Alternatively, the problem can be reduced to the equivalence of weighted automata over a field~\cite[Section B.3]{Bojanczyk2026}.}

\paragraph*{Our Contribution.} The contribution of this paper is \cref{thm:main}, which presents a complete, finitary equational axiomatisation of the rational functions. In order to have such an axiomatisation, we need to develop an algebraic syntax for rational functions, the \emph{rational terms} in which new \kl{rational terms} are constructed from simpler ones by applying certain operations (similarly to regular expressions, which are constructed using sum, concatenation and Kleene star). We therefore begin by defining such an algebraic syntax, although this can hardly be called an original contribution, since our syntax is obtained by  isolating a fragment  of a more expressive system from~\cite{BojanczykDaviaudKrishna2018}. The idea behind our system is that it describes rational functions as morphisms in a certain category, in which the objects are constructed from a singleton \kl{type} $\onetype$ by applying constructors for products, coproducts and lists. A typical morphism in this category would have a type like
\begin{align*}
(\onetype + \onetype)^* \to \onetype^*,
\end{align*}
which describes rational functions that input strings over a \kl{binary alphabet} and output strings over a unary alphabet. The system is defined by identifying certain \kl{atomic terms}, such as the operations of the list monad
\begin{align*}
A^{**} \xrightarrow \concatterm A^* \qquad \text{and} \qquad A \xrightarrow \pureterm A^*,
\end{align*}
and then closing them under certain combinators, such as the composition $f\comp{}g$. The atomic terms and combinators are defined so that they represent exactly the rational functions. The equational axiomatisation consists of axioms such as
\begin{equation}
    \label{eq:intro-examples-of-equations}
\myunderbrace{(f\comp{}g)\comp{}h \equiv f\comp(g\comp{}h)}{associativity of composition}
\qquad \text{or} \qquad 
\myunderbrace{\pureterm\comp\concatterm \equiv \idterm}{one of the axioms  of the list monad}.
\end{equation}

\Cref{thm:main} states that this axiomatisation is finitary and complete. This means that if two terms describe the same rational function, then their equality can be derived from a finite set of \emph{uniform families of axioms}. The axiomatisation is not finite in the literal sense, since the algebra has infinitely many sorts, which are parameterised by types: e.g.~the two axioms in~\eqref{eq:intro-examples-of-equations} are infinite families parameterised by the types of the terms involved. Moreover, unlike Kozen's axiomatisation of Kleene algebra \cite{Kozen1994}, ours uses no implications, only equalities.

\paragraph*{Canonical Models.} Our completeness result requires a substantial proof. One of the central difficulties is related to canonical models. In Kozen's completeness proof for Kleene algebra, the strategy is to prove that every  regular expression is equivalent to one  that describes a minimal automaton~\cite[Lemma 5.4]{Kozen1994}. Since the minimal automaton is unique up to isomorphism, every two regular expressions that describe the same language can be proved equivalent. Unfortunately, this strategy does not work for rational functions, due to the lack of a truly satisfactory canonical model (see~\cite[Section 4]{Filiot2018}). To overcome this difficulty, we use a strategy which is described in \cref{sec:rational-functions}, and which may be of interest outside the context of equational axiomatisations: although rational functions might not have canonical models, such models can be recovered if the input domain is appropriately annotated. Furthermore, the annotations are coherent: two rational functions of the same type can use the same annotation, and thus their canonical models (after annotation) can be meaningfully compared. This is, however, only one of many difficulties, and the others will be described as they arise.

\paragraph*{Structure of the Paper.} In \cref{sec:rational-functions}, we define the rational functions and present a simplified version of our equational axiomatisation. In \cref{sec:syntax}, we define rational terms, an algebraic syntax for rational functions based on~\cite{BojanczykDaviaudKrishna2018}. In \cref{sec:equational-theory} we state the axioms of our equational axiomatisation, and in Sections~\ref{sec:proof-overview}--\ref{sec:completeness-for-all-rational-terms} we prove its completeness. Technical proofs can be found in \appendixplace\iffullversion\else{} of this paper, which can be found on arXiv~\cite{fullversion}\fi.